\chapter{Hardware Comparison}\label{ap:hardware_comparison}

\section{Candidate Boards}

\subsection{NXP S32N74 + ARA-2}

The S32N74\footnote{The S32N74 is in preproduction, so it is not possible to acquire a development board for it.} \cite{S32N7SuperIntegrationProcessor} is a central compute processor designed for software-defined vehicles, enabling the integration of core automotive functions such as safety, scalability, and \gls{AI}-powered processing. 

To further enhance its capabilities, the ARA-2 Module \cite{ARA2} can be used. This module allows embedded systems to be easily upgraded to support Generative \gls{AI} workloads through high-efficiency \gls{eTOPS} performance in a compact form factor.

\subsection{i.MX 95 Evaluation Kit}

The i.MX 95 \cite{IMX95} applications processor family provides a secure and safety-enabled platform. It is specifically designed for automotive, industrial, and \gls{IoT} edge applications, featuring dedicated \gls{ML} acceleration to handle complex vision and data processing tasks.

The i.MX 95 \gls{EVK}\footnote{It's risky to acquire the i.MX 95 \gls{EVK} due to its preproduction status, its specifications can change.} \cite{IMX95EVKQSG} is built upon a modular architecture consisting of an i.MX 95 applications processor and a base board. A key highlight of this platform is the integrated eIQ Neutron \gls{NPU}, which provides dedicated \gls{ML} acceleration for vision and data processing tasks.

\subsection{i.MX 93 Evaluation Kit}

The i.MX 93 \cite{IMX93} family focuses on power-efficient processing for the edge. It provides specialized \gls{ML} acceleration, making it an ideal choice for entry-level smart devices and industrial applications requiring a balance between performance and energy consumption.

The i.MX 93 \gls{EVK} \cite{IMX93Evaluation} offers a platform for developing power-efficient edge devices. The heart of the system is the i.MX 93 processor, which is optimized for energy-efficient \gls{ML} workloads, such as keyword spotting, image recognition, and sensor fusion.

\subsection{NVIDIA Jetson Orin Nano Developer Kit}
    
The NVIDIA Jetson Orin Nano Developer Kit \cite{NVIDIAJetsonAGX} belongs to NVIDIA's latest generation of embedded \gls{AI} computers. As the most affordable option in the Jetson Orin series, it features an 8GB Jetson Orin Nano \gls{AI} module. 

This platform leverages NVIDIA's unified memory architecture, allowing the integrated GPU and CPU to share system resources dynamically. This significantly reduces data transfer overheads during deep learning inference, making it highly attractive for running dense convolutional neural networks and complex computer vision pipelines directly at the edge.

\subsection{Raspberry Pi 5 + Hailo-8}
    
The \gls{RPi} 5 (8GB) \cite{ltdBuyRaspberryPib} constitutes the latest generation of affordable microcomputers developed by the Raspberry Pi Foundation, incorporating all essential computing components.
    
For image classification tasks, the Hailo-8 module \cite{AIAcceleratorHailo8} delivers high-efficiency \gls{AI} acceleration. The \textbf{Hailo-8L} variant \cite{Hailo8LEntryLevelAIa} serves as a cost-effective alternative, offering reduced AI processing capabilities.
    
To facilitate seamless integration and deployment, the \gls{AI} HAT+ \cite{AIHATsRaspberry} has been introduced, available in two versions: one featuring 26 \gls{TOPS} (Hailo-8) and the other offering 13 \gls{TOPS} (Hailo-8L).

\subsection{Raspberry Pi 5 + \texorpdfstring{\acrshort{RPi}}{RPi} \texorpdfstring{\acrshort{AI}}{AI} Camera}

The Raspberry Pi \gls{AI} Camera \cite{AICameraRaspberry} integrates the Sony IMX500 Intelligent Vision Sensor. This hardware allows for the development of sophisticated vision-based \gls{AI} applications by processing neural network models directly on the sensor, reducing the computational load on the main processor.

\section{Comparison}

\subsection{Technical Specifications}

The core computing specifications, peripheral support, and physical constraints of the primary available platforms are compiled below to assess their general system capabilities.

\clearpage

\begin{table}[H]
    \centering
    \caption{Technical Specifications}
    \label{tab:specs}
    \renewcommand{\arraystretch}{1.1}
    %\footnotesize
    \rotatebox{90}{
    \begin{tabular}{p{0.15\textwidth}
                    p{0.28\textwidth}
                    p{0.28\textwidth}
                    p{0.28\textwidth}
                    p{0.28\textwidth}}
        \toprule
        \textbf{Specification} 
                & \textbf{NXP i.MX95 \gls{EVK}} 
                & \textbf{NXP i.MX93 \gls{EVK}} 
                & \textbf{NVIDIA Jetson Orin Nano} 
                & \textbf{\acrshort{RPi} 5} \\
        \toprule

        % GPU
        \textbf{\acrshort{GPU}} 
                & OpenGL ES 3.2, Vulkan 1.2, OpenCL 3.0
                & N/A
                & 1024-core NVIDIA Ampere \acrshort{GPU} 
                & VideoCore VII \acrshort{GPU} \\
        \midrule

        % CPU
        \textbf{\acrshort{CPU}} 
                & 6-core \acrshort{ARM} Cortex-A55 + 1x Cortex-M33 + 1x Cortex-M7
                & 2-core \acrshort{ARM} Cortex-A55 + 1x Cortex-M33 
                & 6-core \acrshort{ARM} Cortex A78AE v8.2 (1.5MB L2 + 4MB L3) 
                & 4-core \acrshort{ARM} Cortex-A76 (512KB/core L2 + 2MB L3) \\
        \midrule
        
        % CPU Max Frequency
        \textbf{\acrshort{CPU} Max Frequency} 
                & 2.0 GHz
                & 1.7 GHz
                & 1.5 GHz 
                & 2.4 GHz \\
        \midrule

        % RAM
        \textbf{\acrshort{RAM}} 
                & 16 GB LPDDR5 (32-bit)
                & 2 GB LPDDR4X 
                & 8 GB 
                & 8 GB \\
        \midrule

        % Storage
        \textbf{Storage} 
                & 64 GB eMMC 5.1, micro\acrshort{SD} slot
                & 16 GB eMMC 5.1, micro\acrshort{SD} slot 
                & \acrshort{SD} Card Slot \& NVMe via M.2 Key M 
                & micro\acrshort{SD} slot (SDR10 mode) \\
        \midrule

        % CSI Camera
        \textbf{\acrshort{CSI} Camera} 
                & 2x \acrshort{MIPI}-\acrshort{CSI} (one muxed with \acrshort{MIPI}-\acrshort{DSI})
                & 1x 2-lane \acrshort{MIPI} \acrshort{CSI} 
                & 2x \acrshort{MIPI} \acrshort{CSI}-2 
                & 2x 4-lane \acrshort{MIPI} for camera/display \\
        \midrule

        % PCIe
        \textbf{\acrshort{PCIe}} 
                & 1x \acrshort{PCIe} x8 slot (audio), 1x \acrshort{PCIe} x4 slot
                & N/A
                & M.2 Key M (x4 \acrshort{PCIe} Gen3), M.2 Key E 
                & M.2 Key M (\acrshort{PCIe} Gen-3.0, 4 lanes) \\
        \midrule

        % USB
        \textbf{\acrshort{USB}} 
                & 1x \acrshort{USB} 2.0 A, 1x \acrshort{USB} 2.0/3.0 C, 1x \acrshort{USB} 2.0 C (Debug)
                & 2x \acrshort{USB} 2.0 A, 1x \acrshort{USB} 2.0 C
                & 1x \acrshort{USB} 3.2 A, 2x \acrshort{USB} 2.0 A, 1x \acrshort{USB} 3.2 C 
                & 2x \acrshort{USB} 3.0, 2x \acrshort{USB} 2.0 \\
        \midrule

        % Wireless
        \textbf{Wireless} 
                & M.2 Key-E (WiFi 6 / \acrshort{BT} 5.3)
                & On-board Module (WiFi 6 /\acrshort{BT} 5.3) 
                & Via \acrshort{GPIO} add-ons 
                & WiFi 2.4/5 GHz 802.11ac, \acrshort{BT} 5.0 \\
        \midrule

        % Power
        \textbf{Power} 
                & 12V / 13.33A (160W AC/DC adapter)
                & 5V / 3A (\gls{USB}-C)
                & 7W--15W 
                & 25W (5V/5A via \acrshort{USB}-C) \\
        \bottomrule
    \end{tabular}
    }
\end{table}

\clearpage

While general computing specs describe basic performance, AI workloads require dedicated hardware accelerators. Table \ref{tab:ai-specs} isolates and contrasts the specialized machine learning performance, software tools, and software platform availability for each system.

\begin{table}[H]
    \centering
    \caption{\acrshort{AI} Specifications}
    \label{tab:ai-specs}
    \renewcommand{\arraystretch}{1.6}
    %\footnotesize
    \begin{tabular}{p{0.2\textwidth} p{0.25\textwidth} p{0.25\textwidth} p{0.2\textwidth}}
        \toprule
        \textbf{Board / Kit} & \textbf{\gls{AI} Performance (\gls{TOPS})} & \textbf{\gls{ML} Frameworks} & \textbf{Supported \acrshort{OS}} \\
        \midrule
        
        \textbf{i.MX 95} & 32 (eIQ Neutron \gls{NPU} \cite{mariasSpotlightEdgeAI2025}) & eIQ \acrshort{ML} (TensorFlow Lite, ONNX, PyTorch \cite{EIQMLSoftware}) & Linux, Android, FreeRTOS \\
        
        \textbf{i.MX 93} & 0.5 (\acrshort{ARM} Ethos U-65 micro\gls{NPU} \cite{NXPIMX93}) & TensorFlow Lite \cite{AN13854NPUMigration} & Linux, FreeRTOS \\
        
        \textbf{NVIDIA Jetson Orin Nano} & 40 & TensorFlow, PyTorch, ONNX, Keras & Linux \\
        
        \textbf{\gls{RPi} 5 + Hailo-8} & 26 & TensorFlow, PyTorch, ONNX, TF Lite & Linux, Windows \\
        
        \textbf{\gls{RPi} 5 + Hailo-8L} & 13 & TensorFlow, PyTorch, ONNX, TF Lite & Linux, Windows \\
        
        \textbf{\gls{RPi} 5 + \acrshort{AI} Cam} & N/A (Sony IMX500 \cite{IMX500DeveloperWorld} \cite{dakicRaspberryPiAI2024}) & TensorFlow Lite, Libcamera & Linux \\
        
        \bottomrule
    \end{tabular}
\end{table}

\subsection{Analysis}

The technical and \gls{AI}-specific data presented in Tables \ref{tab:specs} and \ref{tab:ai-specs} highlight significant architectural differences for deployment environments. The following points summarize the most critical trade-offs:

\begin{itemize}
    \item \textbf{Performance vs. Flexibility}: The \textbf{NVIDIA Jetson Orin Nano} leads in raw \gls{AI} throughput with 40 \gls{TOPS}. Its architecture allows for the highest flexibility when deploying diverse neural network architectures. However, the \textbf{i.MX 95}, with its 32 \gls{TOPS} via the dedicated eIQ Neutron \gls{NPU}, offers a highly efficient alternative for vision-specific tasks, often achieving lower latency in deterministic industrial environments.
    
    \item \textbf{Heterogeneous Computing and Safety}: A key differentiator for the NXP platforms (\textbf{i.MX 95} and \textbf{i.MX 93}) is the inclusion of real-time domains (\gls{ARM} Cortex-M cores). This allows these boards to handle safety-critical tasks, a feature notably absent in the \textbf{Raspberry Pi 5} and less specialized in the \textbf{Jetson} series.
    
    \item \textbf{Modular \gls{AI} Acceleration}: The combination of \textbf{Raspberry Pi 5} and \textbf{Hailo-8} provides a mid-range performance tier (13-26 \gls{TOPS}) that is highly cost-effective for prototyping. While the Hailo modules are power-efficient, the overall system lacks the industrial-level and wide-temperature support found in the NXP and NVIDIA variants.
    
    \item \textbf{Power Efficiency at the Edge}: The \textbf{i.MX 93} and the \textbf{Raspberry Pi \gls{AI} Camera} target "Always-on" applications. These solutions prioritize thermal management and minimal power consumption over high data rates, making them suitable for battery-powered or passively cooled edge devices.
    
    \item \textbf{Software Ecosystems}: NVIDIA's resources remains the most mature for rapid \gls{AI} development. In contrast, NXP's eIQ toolkit focuses on deployment stability and \gls{LTS}, which is critical for automotive and industrial life cycles. The Raspberry Pi ecosystem, while vast, is primarily community-driven and lacks the formal certifications required for critical infrastructure.
\end{itemize}

\section{Final Recommendation}

To determine the ideal deployment targets for this framework, candidate platforms must be matched against practical constraints such as availability, software maturity, and target use case demands. While preproduction models offer high potential, physical prototypes require immediate hardware access and stable \glspl{SDK}. 

The following summary outlines the primary application for each platform based on its technical strengths:

\begin{itemize}
    \item \textbf{NXP i.MX 95}: Recommended for \textbf{industrial and automotive vision systems}. It balances high \gls{AI} performance (32 \gls{TOPS}) with essential safety features. It is the preferred professional platform for deterministic environments where reliability and safety are as critical as processing power.
    
    \item \textbf{NXP i.MX 93}: Specifically designed for \textbf{entry-level, power-efficient smart devices}. Its specialized micro\gls{NPU} (0.5 \gls{TOPS}) is optimized for "always-on" edge tasks such as keyword spotting or simple image recognition. It is the best fit for industrial applications where thermal management and low energy consumption are the primary constraints.
    
    \item \textbf{NVIDIA Jetson Orin Nano}: Best suited for \textbf{high-performance \gls{AI} research and advanced computer vision}. With the highest raw throughput of 40 \gls{TOPS} and a mature software ecosystem, it is the ideal choice for projects requiring complex neural networks or multiple simultaneous video streams in non-safety-critical environments.

    \item \textbf{Raspberry Pi 5 + Hailo-8}: The optimal choice for \textbf{advanced modular prototyping}. Providing 26 \gls{TOPS}, this configuration allows developers to run high-frame-rate inference or multiple models simultaneously on a familiar Linux environment without the cost of high-end industrial \glspl{SoC}.

    \item \textbf{Raspberry Pi 5 + Hailo-8L}: Targeted at \textbf{entry-to-mid-range \gls{AI} applications}. Offering 13 \gls{TOPS}, it serves as a cost-effective entry point for developers who need hardware acceleration significantly beyond a standard \gls{CPU} but do not require the full throughput of the standard Hailo-8 module.

    \item \textbf{Raspberry Pi \gls{AI} Camera}: Targeted at \textbf{integrated, low-power vision sensing}. By processing neural models directly on the Sony IMX500 sensor, it reduces the computational load on the main \gls{CPU}. This makes it suitable for compact, battery-powered edge devices focused on specific vision-based tasks.
\end{itemize}

Ultimately, to validate the deployment automation and cross-platform portability of the unified framework designed in this thesis, a heterogeneous subset of these available boards will be selected. This ensures that the system's runtime abstraction can be tested against drastically different execution backends. A concise mapping of these targets is illustrated in Table \ref{tab:recommendations}.

\begin{table}[H]
    \centering
    \caption{Recommended Use Cases}
    \label{tab:recommendations}
    \renewcommand{\arraystretch}{1.5}
    \begin{tabular}{p{0.3\textwidth} p{0.6\textwidth}}
        \toprule
        \textbf{Platform} & \textbf{Primary Use Case} \\
        \midrule
        \textbf{i.MX 95} & Industrial-grade vision with real-time safety requirements. \\
        \textbf{i.MX 93} & Low-power, always-on edge sensing and control. \\
        \textbf{NVIDIA Jetson Orin Nano} & High-performance \gls{AI} research and complex vision. \\
        \textbf{\gls{RPi} 5 + Hailo-8} & High-speed modular prototyping and advanced \gls{AI} tasks. \\
        \textbf{\gls{RPi} 5 + Hailo-8L} & Cost-effective prototyping for mid-range \gls{AI} workloads. \\
        \textbf{\gls{RPi} 5 + \gls{AI} Camera} & Compact, sensor-level \gls{AI} vision processing. \\
        \bottomrule
    \end{tabular}
\end{table}